% Why: starting point for a paper made with the academic-writing skill; the venue's own class and template replace this one when they exist.
% Build (from the project folder): uv run --with matplotlib python make_figures.py && cd paper && latexmk main.tex
% Tagged, accessible PDF: uncomment the next line with a LaTeX release from 2025 or later, if the venue allows it.
% \DocumentMetadata{lang=en, pdfstandard=ua-2, tagging=on}
\documentclass[11pt,a4paper]{article}
\usepackage[brazilian,english]{babel} % the last language is the paper's main one
\usepackage{microtype}
\usepackage{graphicx}
\usepackage{booktabs}
\usepackage{siunitx}
\sisetup{mode=match} % numbers in the text font
\usepackage[backend=biber,style=authoryear]{biblatex} % citation style: the venue's, or numeric/authoryear by field
\addbibresource{references.bib}
\usepackage[hidelinks]{hyperref}
\usepackage{cleveref} % load after hyperref
\input{../numbers.tex} % written by make_figures.py

\title{<Title that states the central contribution>}
\author{<Name>\\<Affiliation>}
\date{\today}

\begin{document}
\maketitle

% Abstract in five moves (Mensh and Kording): context, gap, what was done, main result with its number, what it means.
\begin{abstract}
<Context in one sentence.> <The gap.> Here we <what was done>.
Over \nRuns{} runs, the new method scored \newMean{} against \baselineMean{} for the baseline, a paired gain of \gainMean{} (range \gainMin{} to \gainMax{}).
<What it means.>
\end{abstract}

\section{Introduction}\label{sec:intro}
% Context, then the gap, then the contributions as claims a reader can check (Peyton Jones). Start with a concrete example.
<Context and a concrete example.> <The gap.> This paper makes the following contributions:
\begin{itemize}
  \item <Claim 1, as listed in claims.md> (\cref{sec:results}).
  \item <Claim 2>.
\end{itemize}

\section{Methods}\label{sec:methods}
% Enough to reproduce: data, cases kept apart for the final check, settings, number of runs, compute, and how the baseline was reproduced and tuned.
<Methods.>

\section{Results}\label{sec:results}
% One subsection per claim; its first sentence states the claim, then the evidence.
\subsection{<Claim 1 as a sentence>}
The new method beats the baseline in every one of the \nRuns{} paired runs (\cref{fig:runs,tab:summary}).
<Why the evidence supports the claim.>

\begin{figure}[t]
  \centering
  \includegraphics{../figures/paper/runs.pdf}
  \caption{The new method scores higher than the baseline in every run. Dots are single runs; bars are means over \nRuns{} runs.}
  \label{fig:runs}
\end{figure}

\begin{table}[t]
  \centering
  \caption{Mean score over \nRuns{} runs and paired gain.}
  \label{tab:summary}
  \begin{tabular}{lr}
    \toprule
    Method & Mean score \\
    \midrule
    Baseline & \baselineMean \\
    New method & \newMean \\
    \midrule
    Paired gain (range) & \gainMean{} (\gainMin{} to \gainMax) \\
    \bottomrule
  \end{tabular}
\end{table}

\section{Discussion}\label{sec:discussion}
% How the gap was filled, the limitations, the relevance to the field (Mensh and Kording).
<What the results mean.> The structure of this template follows \textcite{mensh2017}.

\subsection{Limitations}
<Where the results do not apply, what was not tested, what could break them.>

\printbibliography
\end{document}
